\subsubsection{Localization and Regularization}\label{subsec:distribution-localization}

Although generalized functions cannot be evaluated at points, they can be restricted to open sets. This is quite straightforward. If $G$ is an open subset of $\Omega$, then $\mathcal{D}(G)$ is naturally embedded in $\mathcal{D}(\Omega)$, and hence every generalized function on $\Omega$ defines a generalized function on $G$ by restriction. Consequently, we shall shall define the following.

\begin{Definition}
\label{def:distribution-vanishing}
We say that $f\in\mathcal{D}^{\prime}(\Omega)$ \textbf{vanishes} on an open set $G\subset\Omega$ if $(f,\phi)=0$ for every $\phi\in\mathcal{D}(G)$. Two distributions are equal on $G$ if their difference vanishes on $G$.
\end{Definition}

It can be shown (cf. Problem~\ref{ex:distribution-local-vanishing}) that if $f$ vanishes locally near every point of $G$, i.e., if every point of $G$ has a neighborhood on which $f$ vanishes, then $f$ vanishes on $G$. An immediate consequence is that if $f$ vanishes on each of a family of open sets, it also vanishes on their union. Hence there is a largest open set $N_f$ on which $f$ vanishes.

\begin{Definition}
\label{def:distribution-support}
The complement of $N_f$ in $\Omega$ is called the \textbf{support} of $f$.
\end{Definition}

\begin{Example}
\label{ex:delta-support}
The support of the delta function is the set $\{0\}$. Although the delta function cannot be evaluated at points, it makes sense to say that it vanishes except at the origin.
\end{Example}

\begin{Remark}
\label{rem:singular-function-regularization}
Functions with nonintegrable singularities are not defined as generalized functions by equation \eqref{eq:regular-distribution-pairing}. However, it is often possible to define a generalized function which agrees with a singular function on any open set that does not contain the singularity. Such a generalized function is called a regularization. For example, a regularization of the function $1/x$ on $\mathbb{R}$ is given by the principal value integral
\begin{equation}\label{eq:principal-value-one-over-x-regularization}
\begin{aligned}
(f,\phi)&=\int_{-\infty}^{-\epsilon}\dfrac{\phi(x)}x\,\mathrm{d}x
+\int_{-\epsilon}^{\epsilon}\dfrac{\phi(x)-\phi(0)}x\,\mathrm{d}x\\
&\qquad+\int_\epsilon^\infty\dfrac{\phi(x)}x\,\mathrm{d}x
\end{aligned}
\end{equation}
(cf. Problem~\ref{ex:principal-value-regularization}).
\end{Remark}
